\documentclass[onecolumn]{article}

\renewcommand*\familydefault{\sfdefault}

\usepackage{amsmath}
\usepackage{bm}
\usepackage{amssymb}
\usepackage{helvet}
\usepackage{graphicx}
\usepackage{epstopdf}

\setlength{\topmargin}{0in}
\setlength{\oddsidemargin}{0in}
\setlength{\headheight}{0.5in}
\setlength{\headsep}{0in}
\setlength{\textheight}{8.5in}
\setlength{\textwidth}{6.5in}

\setlength{\parindent}{0.2in}
\setlength{\parskip}{0.1in}

\begin{document}

\begin{center}

{\Huge Finding The Time Constant of a Circuit}\vskip 20pt

{\huge \em Devesh Parekh}

\end{center}

\section*{Introduction}

The purpose of this lab is to find the time constant, $\tau$, in an RC circuit.This was done by connecting a function generator to a capacitor, then to a resistor, which was then connected to an oscilloscope. We then adjusted the oscilloscope to display a square wave. We then adjusted the vertical and horizontal knobs to align the graph displayed on the oscilloscope. The y-axis of the graph represents the voltage, and the x-axis represents the time in seconds. From the graph displayed on the oscilloscope, we gathered 10 data points for voltage and time. The final result of the experiment is as follows: $\tau = 1.4163 \pm 0.11729\,s$.

The fundamental equation for a discharging RC circuit was used for this lab:

\begin{equation}
\tau = RC
\end{equation}

In this equation, the time constant $\tau$ is equal to the resistance $R$ times the capacitance $C$.The above equation can then be used to derive the discharging capacitor equation below:

\begin{equation}
V = V_0e^{-t/RC}
\end{equation}

In this equation, the voltage $V$ is equal to the initial voltage $V_0$ times Euler's number $e$ raised to the power of $-t$ divided by the resistance $R$ times the capacitance $C$.

\begin{equation}
S = -\frac{1}{\tau}
\end{equation}

Equation (2) is in the form of the equation of a line, where the slope is given by $S=-\frac{1}{\tau}$.

\section*{Methods}

We started by connecting the function generator to a capacitor. The capacitor was then connected to a resistor.Finally, the resistor was connected to an oscilloscope. We set the function generator to a square wave. The oscilloscope then displayed a wave on an axis. The axis could be adjusted on the oscilloscope using knobs, which adjusted the graph horizontally and vertically. After we adjusted the graph appropriately, we defined the axes as a class. The x-axis represented time, and the y-axis represented voltage. We then wrote down the values on the axes for voltage and time. This process was repeated until we had 10 data points for both $V$ and $T$.

The value of the slope is determined from a least squares fit to time $T$ plotted as a natural log function of voltage $V$. The value of $\tau$ is calculated using $V$ and $T$. The uncertainty in $\tau$ is calculated by propagating the uncertainty for $S$ in the calculation for $\tau$. Unless otherwise noted, all uncertainties are reported using $t$-statistics for 95\% confidence intervals and confidence bands. All calculations are performed in MATLAB\footnote{The MathWorks, http://www.mathworks.com}.

\section*{Data}

The data are shown in Table 1. The voltage and time for this lab were measured using the oscilloscope.The uncertainty of the function generator is $\pm0.5~\mu$s, and the uncertainty of the oscilloscope is $\pm0.05$ volts.

\begin{table}[htbp!]
\centering
\caption{Data}
\begin{tabular}{cccc}
\hline\hline
$I$ & $r$ \\
($amps$) & ($cm$) \\
\hline
1.1 & 4 \\
1.2 & 4.15 \\
1.3 & 4.5 \\
1.4 & 4.6 \\
1.6 & 4.85 \\
1.7 & 5 \\
1.8 & 5.15 \\
1.9 & 5.5 \\
2.0 & 5.75 \\
\hline
\end{tabular}
\end{table}

\section*{Results}

The results for the fitted model are shown in Figure 1.The slope of the line and the value of the time constant are $\tau = 1.4163 \pm 0.11729\,s$.

\begin{figure*}[t!]
\includegraphics[width=1\textwidth]{fig_t.eps}
\caption{Least-squares fit to $\ln(V)$ versus time. The data are shown as blue points.The fitted model is shown as a red solid line. The 95\% confidence bands are shown as green dashed lines.}
\end{figure*}

The results for the fitted model are shown in the image above.The results show that our experiment is not comparable to the expected value of the RC circuit. If we input the values for the resistor ($10~k\Omega$) and the capacitor ($68~nF$) into the equation for the time constant, $\tau = RC$, we get a value of $6.8 \times 10^{-4}\,s$, which is much less than what we calculated. This experiment shows that the real-world value may be different from the theoretical value due to other factors that we are not aware of or did not account for. My hunch is that it has something to do with the impedance of the oscilloscope.

\end{document}